\documentclass[11pt]{article}
\usepackage[margin=1in]{geometry}
\usepackage{amsmath,amssymb,mathtools,bm}
\usepackage{booktabs}
\usepackage{array}
\usepackage{hyperref}

\title{Why the Bounded-\texorpdfstring{$u_{\mathrm{ref}}$}{u\_ref} Run Produces a Very Small \texorpdfstring{$x_s$}{x\_s}}
\author{}
\date{\today}

\begin{document}
\maketitle

\section{Question}

The observed behavior is:
\begin{itemize}
\item in the \emph{unbounded} frozen-$\hat d$ run, the online target state $x_s$ is large,
\item in the \emph{bounded with $u_{\mathrm{ref}}$} run, the online target state $x_s$ stays close to zero,
\item this raises the question of whether the steady-state equations were implemented incorrectly, or whether $x_s$ is simply very dependent on $u_s$.
\end{itemize}

This note answers that question using the actual matrices and the latest two runs:
\begin{itemize}
\item bounded with $u_{\mathrm{ref}}$: \texttt{Data/debug\_exports/mpc\_first\_step\_replacement\_bounded\_frozen\_dhat/20260406\_232910}
\item unbounded: \texttt{Data/debug\_exports/mpc\_first\_step\_replacement\_unbounded\_frozen\_dhat/20260406\_232953}
\end{itemize}

\section{Bottom Line}

The implementation is internally consistent. The evidence strongly supports:
\begin{enumerate}
\item the code is using the correct steady-state algebra,
\item $x_s$ is exactly determined from $u_s$ through
\begin{align}
x_s = (I-A)^{-1}Bu_s,
\end{align}
\item the very small bounded target state does \emph{not} come from a wrong $x_s$ formula,
\item it comes from the fact that the bounded-\(u_{\mathrm{ref}}\) optimization is choosing a very small $u_s$.
\end{enumerate}

The key structural point is:
\begin{align}
\text{small bounded } x_s
\quad\Longleftarrow\quad
\text{small bounded } u_s
\quad\Longleftarrow\quad
\lambda_u \|u_s-u_{\mathrm{ref}}\|_2^2 \text{ dominates } \|Gu_s-r_k\|_2^2.
\end{align}

\section{Steady-State Equations Used in the Code}

For the frozen-$\hat d$ target, the code solves
\begin{align}
(I-A)x_s - Bu_s &= 0, \\
Cx_s &= r_k,
\end{align}
where
\begin{align}
r_k = y_{sp,k} - \hat d_k.
\end{align}

If $(I-A)$ is invertible, then
\begin{align}
x_s &= (I-A)^{-1}Bu_s, \\
Gu_s &= r_k, \qquad G := C(I-A)^{-1}B.
\end{align}

Define
\begin{align}
W := (I-A)^{-1}B.
\end{align}
Then
\begin{align}
x_s = Wu_s,
\qquad
y_s - d_s = Cx_s = CWu_s = Gu_s.
\end{align}

This means:
\begin{itemize}
\item the output-side steady-state relation is controlled by $G$,
\item the state-side steady-state relation is controlled by $W$.
\end{itemize}

\section{Actual Identified Matrices}

For the current notebook setup, the identified unaugmented model is
\begin{align}
A =
\begin{bmatrix}
0.946394 & 0 & 0 & 0 & 0 & 0 & 0 \\
0 & 0.936731 & 0 & 0 & 0 & 0 & 0 \\
0 & 0 & 0.885846 & 0 & 0 & 0 & 0 \\
0 & 0 & 0 & 0.927830 & 0 & 0 & 0 \\
0 & 0 & 0 & 0 & 0 & 1 & 0 \\
0 & 0 & 0 & 0 & 0 & 0 & 1 \\
0.010342 & 0 & 0 & 0 & 0 & 0 & 0
\end{bmatrix},
\end{align}
\begin{align}
B =
\begin{bmatrix}
7.601178\times 10^{-3} & 0 \\
1.512532\times 10^{-2} & 0 \\
0 & 1.471516\times 10^{-2} \\
0 & 3.010827\times 10^{-2} \\
0 & 0 \\
0 & 0 \\
3.267612\times 10^{-5} & 0
\end{bmatrix},
\end{align}
\begin{align}
C =
\begin{bmatrix}
0 & 0 & 0.029531 & 0 & 1 & 0 & 0 \\
0 & -0.050171 & 0 & 0.046562 & 0 & 0 & 0
\end{bmatrix}.
\end{align}

The matrix $(I-A)$ is not near singular:
\begin{align}
\kappa(I-A) \approx 33.62.
\end{align}
So there is no evidence that the $x_s$ calculation is blowing up because $(I-A)^{-1}$ is numerically invalid.

\section{State Map and Output Map}

Using the actual model,
\begin{align}
W = (I-A)^{-1}B =
\begin{bmatrix}
0.141797 & 0 \\
0.239063 & 0 \\
0 & 0.128906 \\
0 & 0.417188 \\
0.001499 & 0 \\
0.001499 & 0 \\
0.001499 & 0
\end{bmatrix},
\end{align}
and
\begin{align}
G = CW =
\begin{bmatrix}
0.001499 & 0.003807 \\
-0.011994 & 0.019425
\end{bmatrix}.
\end{align}

Its inverse is
\begin{align}
G^{-1} =
\begin{bmatrix}
259.770673 & -50.907293 \\
160.396561 & 20.047199
\end{bmatrix}.
\end{align}

\subsection{Interpretation}

These numbers immediately explain the behavior:
\begin{itemize}
\item $W$ has moderate gain. The singular values of $W$ are
\begin{align}
\sigma(W) = [0.436649,\; 0.277964].
\end{align}
So $x_s$ is \emph{not} exploding relative to $u_s$.
\item The real amplification is in $G^{-1}$. A modest output-side right-hand side $r_k$ can require a very large exact $u_s$.
\end{itemize}

So the main story is not
\begin{quote}
``$x_s$ is hypersensitive to $u_s$.''
\end{quote}
It is instead
\begin{quote}
``the exact output-tracking target requires a huge $u_s$, and then $x_s$ follows that $u_s$ linearly through $W$.'' 
\end{quote}

\section{Run-to-Run Comparison}

Comparing the latest two runs gives:
\begin{center}
\begin{tabular}{lcc}
\toprule
Metric & bounded with $u_{\mathrm{ref}}$ & unbounded \\
\midrule
mean $\|u_s\|_2$ & 5.0394 & 449.3150 \\
mean $\|x_s\|_2$ & 1.7323 & 149.1337 \\
ratio to unbounded & 0.0112 & 1.0000 \\
ratio of $\|x_s\|_2$ & 0.0116 & 1.0000 \\
\bottomrule
\end{tabular}
\end{center}

The crucial observation is that the shrinkage ratios are almost the same:
\begin{align}
\frac{\mathbb{E}\|x_s^{\mathrm{bounded}}\|}{\mathbb{E}\|x_s^{\mathrm{unbounded}}\|} &\approx 0.0116, \\
\frac{\mathbb{E}\|u_s^{\mathrm{bounded}}\|}{\mathbb{E}\|u_s^{\mathrm{unbounded}}\|} &\approx 0.0112.
\end{align}

That is exactly what one expects from
\begin{align}
x_s = Wu_s
\end{align}
with a moderate-gain matrix $W$.

\section{Direct Consistency Check}

For both runs, the saved online target satisfies the steady-state state equation numerically:
\begin{align}
x_s \approx Wu_s.
\end{align}
The measured reconstruction residuals were:
\begin{align}
\max_k \|x_s^{\mathrm{bounded}}(k) - W u_s^{\mathrm{bounded}}(k)\|_2
&= 5.44\times 10^{-16}, \\
\max_k \|x_s^{\mathrm{unbounded}}(k) - W u_s^{\mathrm{unbounded}}(k)\|_2
&= 8.16\times 10^{-14}.
\end{align}

Also, the run-to-run difference satisfies
\begin{align}
x_s^{\mathrm{unbounded}} - x_s^{\mathrm{bounded}}
=
W\left(u_s^{\mathrm{unbounded}} - u_s^{\mathrm{bounded}}\right)
\end{align}
to machine precision:
\begin{align}
\max_k
\left\|
\left(x_s^{\mathrm{unbounded}} - x_s^{\mathrm{bounded}}\right)
- W\left(u_s^{\mathrm{unbounded}} - u_s^{\mathrm{bounded}}\right)
\right\|_2
=
1.15\times 10^{-13}.
\end{align}

This is very strong evidence that the implementation is algebraically correct.

\section{Why the Bounded-\texorpdfstring{$u_{\mathrm{ref}}$}{u\_ref} Target Is So Small}

The bounded-with-anchor objective is
\begin{align}
\min_{u_s}
\quad
\|Gu_s-r_k\|_2^2 + \lambda_u \|u_s-u_{\mathrm{ref}}\|_2^2
\qquad
\text{subject to }
u_{\min} \le u_s \le u_{\max}.
\end{align}

In the current run,
\begin{align}
\lambda_u = 1.0,
\qquad
u_{\mathrm{ref}} = u_{k-1}.
\end{align}

The Hessian of the output-tracking term is
\begin{align}
G^\top G =
\begin{bmatrix}
1.4610\times 10^{-4} & -2.2728\times 10^{-4} \\
-2.2728\times 10^{-4} & 3.9182\times 10^{-4}
\end{bmatrix},
\end{align}
with eigenvalues
\begin{align}
\lambda(G^\top G) = [1.0604\times 10^{-5},\; 5.2732\times 10^{-4}].
\end{align}

This is the crucial scale comparison:
\begin{align}
G^\top G \sim 10^{-5}\text{ to }10^{-4},
\qquad
\lambda_u I \sim 1.
\end{align}

So the anchor term is about \(10^3\) to \(10^5\) times stronger than the output-tracking curvature.

\subsection{Inactive Bounds in the Actual Run}

In the bounded-with-$u_{\mathrm{ref}}$ run, the target never actually sat on the input bounds:
\begin{align}
\text{fraction active at lower bounds} &= [0,\;0], \\
\text{fraction active at upper bounds} &= [0,\;0].
\end{align}

So in this particular run, the small bounded target is \emph{not} mainly caused by clipping on the box. It is mainly caused by the $u_{\mathrm{ref}}$ penalty.

\subsection{Step-0 Proof}

At step \(k=0\), the exact unbounded target was
\begin{align}
u_s^{\mathrm{unbounded}} =
\begin{bmatrix}
689.123462 \\
451.554842
\end{bmatrix},
\end{align}
which corresponds to
\begin{align}
r_0 = Gu_s^{\mathrm{unbounded}} =
\begin{bmatrix}
2.752065 \\
0.506106
\end{bmatrix}.
\end{align}

Because the bounded solution at this step is interior, the regularized unconstrained minimizer is
\begin{align}
u_s^{\mathrm{reg}}
=
(G^\top G + \lambda_u I)^{-1}(G^\top r_0 + \lambda_u u_{\mathrm{ref}}).
\end{align}
With \(u_{\mathrm{ref}} = 0\) and \(\lambda_u = 1\), this gives
\begin{align}
u_s^{\mathrm{reg}}
=
\begin{bmatrix}
-0.001940 \\
0.020300
\end{bmatrix},
\end{align}
which matches the actual bounded online target:
\begin{align}
u_s^{\mathrm{bounded}}
=
\begin{bmatrix}
-0.001939 \\
0.020298
\end{bmatrix}.
\end{align}

This is the cleanest evidence in the whole analysis:
\begin{quote}
the tiny bounded target is exactly what the weighted optimization asks for, even before the box constraints matter.
\end{quote}

\section{Why \texorpdfstring{$x_s$}{x\_s} Then Becomes Tiny}

Once the optimization chooses a tiny \(u_s\), the state target is forced to be tiny because
\begin{align}
x_s = Wu_s.
\end{align}
At step \(k=0\),
\begin{align}
x_s^{\mathrm{bounded}}
&=
\begin{bmatrix}
-2.75\times 10^{-4} \\
-4.64\times 10^{-4} \\
2.62\times 10^{-3} \\
8.47\times 10^{-3} \\
-2.91\times 10^{-6} \\
-2.91\times 10^{-6} \\
-2.91\times 10^{-6}
\end{bmatrix}, \\
x_s^{\mathrm{unbounded}}
&=
\begin{bmatrix}
97.715553 \\
164.743578 \\
58.208241 \\
188.383036 \\
1.033049 \\
1.033049 \\
1.033049
\end{bmatrix}.
\end{align}

This is not because \(x_s\) was coded incorrectly. It is because the chosen \(u_s\) changed by orders of magnitude, and \(x_s\) follows it linearly.

\section{Answer}

The correct answer is:
\begin{enumerate}
\item Yes, \(x_s\) is directly determined by \(u_s\) through the exact relation
\begin{align}
x_s = (I-A)^{-1}Bu_s.
\end{align}
\item No, the evidence does \emph{not} suggest the code was written incorrectly.
\item The actual source of the discrepancy is the optimization scaling:
\begin{itemize}
\item \(G^{-1}\) is large, so the exact output-tracking target wants a huge \(u_s\),
\item \(W\) is moderate, so \(x_s\) then becomes proportionally large,
\item with \(\lambda_u = 1\), the \(u_{\mathrm{ref}}\) penalty dominates \(G^\top G\),
\item therefore the bounded-with-\(u_{\mathrm{ref}}\) optimization picks a tiny \(u_s\),
\item and then \(x_s = Wu_s\) is also tiny.
\end{itemize}
\end{enumerate}

\section{Practical Consequence}

If the goal is for the bounded target to remain closer to the exact frozen-$\hat d$ equilibrium, then the main tuning issue is not the \(x_s\) formula. The main tuning issue is the relative scale of
\begin{align}
\|Gu-r_k\|_2^2
\quad\text{versus}\quad
\lambda_u \|u-u_{\mathrm{ref}}\|_2^2.
\end{align}

With the present model,
\begin{align}
\lambda_u = 1
\end{align}
is extremely large relative to the natural scale of \(G^\top G\). That is why the bounded-with-anchor target stays very close to \(u_{\mathrm{ref}}\) and therefore produces a very small \(x_s\).

\end{document}
